\chapter{记号表 Notation}

本附录统一整理全书符号约定，并显式标出各讲出现的局部重载。目标不是强行把所有 lecture 的原始记号改成一套，而是让读者能判断：\textbf{同一个符号在全书中默认表示什么，何时被局部特化。}

\section{全书记号治理规则}

\begin{itemize}
\item 默认优先复用通用对象类：输入用 $x$，输出用 $y$，策略模型用 $\pi_\theta$，参考策略用 $\pi_{\mathrm{ref}}$，预算用 $B$，轨迹用 $\tau$。
\item 若同一符号在不同章节表示同一类对象的特化，应通过下标或文字限定说明，而不是重新发明完全无关的新符号。
\item 若符号发生真正重载，正文必须在第一次出现时显式写出本地语义。
\item $V$ 仅在明确表示 verifier 或 verification function 时使用；value estimation 优先写成 $v_\phi(s)$ 或 $Q(s,a)$。
\item $a$ 在全书中默认表示 action；在 theorem proving 章节，tactic 被视为 proof environment 中的特殊 action。
\item $s$ 在全书中默认表示 state；在证明章节中是 proof state，在安全章节中可表示 agent context state。
\end{itemize}

\section{基础输入输出与偏好优化记号}

\begin{description}
\item[$x$] 输入 query、prompt 或 task instance 的默认记号。除非章节另行声明，全书都按此理解。

\item[$y$] 输出 answer、response、completion 或 formal candidate 的默认记号。

\item[$\hat{y}$] 聚合或最终预测得到的答案，常见于 self-consistency、voting 或 verifier-backed aggregation。

\item[$y_w, y_l$] preference pair 中的 winner / loser response。主要出现在 DPO、IRPO 和相关偏好优化章节。

\item[$\pi_\theta$] 当前待训练或待部署的策略模型。

\item[$\pi_{\mathrm{ref}}$] 参考策略模型，用于约束偏离幅度或定义 preference objective。

\item[$\mathcal{D}_t$] 第 $t$ 轮训练或迭代中的数据集。

\item[$\lambda$] 平衡多项目标的系数；在本书中主要出现在 IRPO 一类混合目标中。
\end{description}

\section{预算、分数、验证与价值函数}

\begin{description}
\item[$B$] 预算（budget）。默认指 inference-time budget 或 interaction budget，不拿来表示 batch size、beam width 之外的其他量。

\item[$s_i$] 对候选对象 $x_i$ 的评分。常用于 prompt optimization 或 candidate ranking。

\item[$r_i$] judge 给第 $i$ 个候选回答的 reward。

\item[$r_\phi$] 由 verifier、reward model 或 judge 参数化给出的评分函数。若具体评分来源重要，应保留下标。

\item[$V(x,y)$] verification function，检查输入 $x$ 与输出 $y$ 是否满足外部正确性条件。主要用于 RLVR 或可验证任务。

\item[$v_\phi(s)$] 状态 $s$ 的 value estimate。用于 web agent tree search、planning 或其它状态价值建模场景。

\item[$Q(s,a)$] 在状态 $s$ 采取动作 $a$ 的 action-value estimate。主要出现在 theorem proving search 或 RL-heavy 语境中。
\end{description}

\section{轨迹、状态与环境交互}

\begin{description}
\item[$\tau$] 轨迹（trajectory）的默认记号。对 coding agents、web agents、GUI agents、多模态 agents，它表示 observation-action-reward style interaction trace；对 theorem proving，它局部特化为 proof trajectory。

\item[$\mathcal{T}$] 候选 trajectories 或 reasoning traces 的集合。

\item[$o_t$] 第 $t$ 步 observation 或 tool output。

\item[$a_t$] 第 $t$ 步 action 或 tool call。在 GUI 章节，它可以细化为坐标化动作；在 theorem proving 章节，它可以细化为 tactic。

\item[$s, s_t$] 当前 state 或第 $t$ 步 state 的默认记号。若上下文是 theorem proving，则指 proof state；若上下文是安全章节，则可指 agent context state。

\item[$\mathcal{S}, \mathcal{A}, \mathcal{O}$] 状态空间、动作空间、观测空间。用于 POMDP 或 environment 抽象。

\item[$\mathcal{M}=(\mathcal{S},\mathcal{A},\mathcal{O},T,R)$] 交互环境的 POMDP 式抽象，其中 $T$ 是 transition dynamics，$R$ 是 reward 或 task signal。
\end{description}

\section{多模态与时间序列专用记号}

\begin{description}
\item[$p^k$] 单步正确率 $p$ 在 $k$ 步任务中的近似成功概率，用于说明长程任务为何随 horizon 快速变难。

\item[$P(\mathrm{task}_i \mid \mathrm{task}_{<i})$] 长程任务分阶段成功概率，用于 Plan-Seq-Learn 式层级分解。

\item[$p_\theta(t_{1:m}, a_{1:n} \mid x)$] CoTA-style multimodal action model 对 thought sequence 与 action sequence 的联合建模。

\item[$a_t=(\mathrm{type}_t,x_t,y_t,\mathrm{arg}_t)$] GUI action 在统一视觉空间中的结构化表示，包括动作类型、坐标与参数。

\item[$z_{1:K}=f_\theta(v_{1:T})$] temporal encoder 把长视频帧序列 $v_{1:T}$ 压缩成较短 token 序列 $z_{1:K}$ 的写法，其中通常有 $K \ll T$。

\item[$\hat{\mathcal{I}}=\{(i,r_i)\}$] 与 frame relevance 或信息密度相关的标注集合。
\end{description}

\section{形式化数学与证明专用记号}

\begin{description}
\item[$x_I$] informal problem statement。

\item[$x_F$] formal problem statement。

\item[$g$] 当前 theorem goal。

\item[$s$] proof state；在 formal proving 章节里，这是 $s$ 的默认特化含义。

\item[$a$] 候选 tactic 或 proof action；在 proving 章节里，这是 action 的局部特化。

\item[$z_t$] 第 $t$ 步 tactic 前生成的 informal thought，用于 thought-augmented proving。

\item[$P_\theta(a \mid s)$] prover model 在 proof state $s$ 下对 tactic $a$ 给出的动作先验。

\item[$f_\phi$] formalizer 或 theorem/proof translator 模型。

\item[$g_\phi$] theorem autoformalizer，一般指把 informal theorem 转成 formal theorem 的模型。

\item[$h_\psi$] proof autoformalizer，一般指把 informal proof 转成 formal proof sketch 或 formal proof 的模型。

\item[$r_\eta(P \mid s)$] premise retrieval 或 premise selection score。

\item[$p$] 候选 premise。

\item[$c_{\mathrm{file}}, c_{\mathrm{repo}}$] 文件级与项目级上下文，用于 long-context theorem proving 或 repository-scale proving。

\item[$T_1 \equiv T_2$] 两个 formal theorem 的语义等价关系。

\item[$\hat{\tau} = \operatorname{Prove}(g \mid s_{\mathrm{local}}, c_{\mathrm{file}}, c_{\mathrm{repo}})$] 在给定局部 state 和长上下文条件下生成 proof trajectory 的抽象写法。

\item[$\mathcal{N}(p)$] 围绕目标问题构造的邻域或变体分布，常用于 AlphaProof 式 problem specialization 讨论。
\end{description}

\section{安全与策略执行专用记号}

\begin{description}
\item[$x$] 在安全章节中可局部重载为 attack seed 或其变体。该语义仅限攻击建模上下文，不覆盖全书默认输入语义。

\item[$\operatorname{ASR}(x)$] attack success rate，衡量攻击样本 $x$ 的成功程度。

\item[$\operatorname{Cov}(x)$] 覆盖度指标，表示攻击、任务或状态分布覆盖程度。

\item[$u$] 用户身份、角色或能力约束。

\item[$P_h$] human-written global policy，即人工编写的高层政策或全局约束。

\item[$P_d$] dynamic policy，即根据上下文动态生成或调整的执行约束。
\end{description}

\section{全书主要重载的裁决}

\begin{itemize}
\item \textbf{$x$}：默认是输入；只有安全章节在攻击建模中把它局部重载为 attack seed。
\item \textbf{$s$}：默认是一般 state；在 proving 章节默认指 proof state；在安全章节如果表示 context state，正文必须显式说清。
\item \textbf{$a$}：默认是一般 action；在 theorem proving 中是 tactic 这一特殊 action。
\item \textbf{$\tau$}：默认是 trajectory；在 theorem proving 中是 proof trajectory，在 abstraction/discovery 章节可泛化为 search trajectory。
\item \textbf{$V$ 与 $v$}：$V(x,y)$ 保留给 verification function；$v_\phi(s)$ 保留给 value estimation。两者不得混写。
\end{itemize}
